% STATUS: ACTIVE -- three-form function set, one page per quantity.
%          Common setting: w_bar_s known, each form's exponent pinned, and every
%          quantity evaluated ALONG b(w_bar) -- the constant that form's law
%          demands at each height.  No fitted constant appears anywhere.
%              Form A   a = 1 - exp(-b*(w-ws))     p = 0
%              Form B   a = 1 - b/((w-ws) + b)     p = 1   (production)
%              Form C   a = 1 - b/(1 + w-ws)       p = 1   (amplitude writing)
%          Five quantities: b(w_bar) , db(w_bar)/dw_bar , a_bar' , da_bar/db , J_b
%          (a_bar is omitted: along b(w_bar) all three collapse onto a_true).
%
%          PAGES 6-8 (2026-08-09) take page 1 seriously and close the loop.
%          Page 1's b(w_bar) is what formB_ws_ref.tex S10 calls b_eff; how far
%          the fixed published anchor sits from it is what decides whether
%          estimating the constant is worth anything, and how flat it is over
%          the envelope is what decides the gain error a frozen constant can
%          reach. Two truth functions are used: the wall-normal c_perp (the
%          production case) and the wall-parallel c_para, the latter SUBSTITUTED
%          into the wall-normal channel through the driver's plant_cperp handle
%          -- same filter, same axis, same trajectory, different truth. It is
%          not "the x/y axes".
%          Figures: test_script/scratch/plot_formB_anchor_vs_target.m (p.6)
%                   test_script/scratch/plot_formB_pages78.m          (p.7-8)
%          Arms:    test_script/scratch/run_formB_BC_two_boundaries.m
%                   (2 truths x 2 writings x free/locked, 12 seeds, 96 runs)
%          Acceptance for the amplitude arm: verify_formB_amp_tie.m (7/7 PASS).
%          NOTATION: b(w_bar) with an argument = the constant the TRUTH demands
%          at each height (formB_ws_ref.tex S10 calls it b_eff); plain b = the one
%          frozen constant the filter carries.  Pages 1-2 use the function, pages
%          3-6 the constant.
%          Results only, no intermediate algebra and no prose (user, 2026-08-08);
%          (w_bar - w_bar_s) written out in full everywhere -- no shorthand.
% ======================================================
% Compile: xelatex formB_amp_functions.tex
% Path:    reference/eq17_analysis/derivation/formB_amp_functions.tex
% Figures: test_script/scratch/plot_formB_form_compare.m
%          -> figures/formB_cmp_{b,dbdw,aprime,dadb,Jb}.png
% Truth:   calc_correction_functions (two-sphere perpendicular series);
%          the body states everything through a_true and a'_true, never c_perp
% ======================================================

\documentclass[12pt,a4paper,fontset=none]{ctexart}

\usepackage[fleqn]{amsmath}
\usepackage{amssymb}
\setlength{\mathindent}{0pt}
\usepackage{geometry}
\geometry{margin=1.8cm}
\usepackage{graphicx}

\setlength{\parindent}{0pt}
\setlength{\parskip}{0.6em}
\linespread{1.05}
\pagestyle{plain}

\newcommand{\gp}{\bigl(\bar{w}-\bar{w}_s\bigr)}
\newcommand{\gq}{\bigl(1+\bar{w}-\bar{w}_s\bigr)}

% one page = results on top, one large figure filling the rest
\newcommand{\fpage}[2]{%
  #1\par\vfill
  \begin{center}\includegraphics[width=0.97\linewidth]{figures/#2}\end{center}%
  \vfill\clearpage}

\begin{document}

% ---------------------------------------------------------------- step 1: b
\fpage{%
\begin{align*}
\text{Form A}\quad \bar{a} &= 1 - e^{-b\gp}       &  p &= 0 \\
\text{Form B}\quad \bar{a} &= 1 - \frac{b}{\gp+b} &  p &= 1 \\
\text{Form C}\quad \bar{a} &= 1 - \frac{b}{1+\bar{w}-\bar{w}_s} &  p &= 1
\end{align*}

\begin{equation*}
\bar{w}_s\ \text{known},
\qquad
\bar{a}^{\prime}_{\mathrm{true}} := \frac{\mathrm{d}\bar{a}_{\mathrm{true}}}{\mathrm{d}\bar{w}}
\end{equation*}

\begin{align*}
b_{\text{A}}(\bar{w}) &= -\frac{\ln\bigl(1-\bar{a}_{\mathrm{true}}\bigr)}{\bar{w}-\bar{w}_s},
&
b_{\text{B}}(\bar{w}) &= \gp\,\frac{1-\bar{a}_{\mathrm{true}}}{\bar{a}_{\mathrm{true}}},
&
b_{\text{C}}(\bar{w}) &= \gq\bigl(1-\bar{a}_{\mathrm{true}}\bigr)
\end{align*}

\begin{equation*}
\bar{w}-\bar{w}_s \in [\,0.1,\ 10\,]
\end{equation*}

\begin{align*}
\bar{w}-\bar{w}_s = 0:&\quad
  \text{A, B}:\ \text{any } b \ ;\qquad \text{C}:\ b = 1 \\
\bar{w}-\bar{w}_s = \infty:&\quad
  \text{A, B, C}:\ \text{any } b \\
\text{elsewhere}:&\quad
  \partial\bar{a}/\partial b \neq 0 \ \Rightarrow\ b(\bar{w}) \ \text{unique}
\end{align*}}{formB_cmp_b.png}

% ------------------------------------------------------------ step 2: db/dw
\fpage{%
\begin{align*}
\frac{\mathrm{d}b_{\text{A}}(\bar{w})}{\mathrm{d}\bar{w}}
  &= \frac{1}{\gp^{2}}
     \left[\frac{\gp\,\bar{a}^{\prime}_{\mathrm{true}}}{1-\bar{a}_{\mathrm{true}}}
     + \ln\bigl(1-\bar{a}_{\mathrm{true}}\bigr)\right] \\[0.5em]
\frac{\mathrm{d}b_{\text{B}}(\bar{w})}{\mathrm{d}\bar{w}}
  &= \frac{1-\bar{a}_{\mathrm{true}}}{\bar{a}_{\mathrm{true}}}
   - \gp\,\frac{\bar{a}^{\prime}_{\mathrm{true}}}{\bar{a}_{\mathrm{true}}^{\,2}} \\[0.5em]
\frac{\mathrm{d}b_{\text{C}}(\bar{w})}{\mathrm{d}\bar{w}}
  &= \bigl(1-\bar{a}_{\mathrm{true}}\bigr) - \gq\,\bar{a}^{\prime}_{\mathrm{true}}
\end{align*}

\begin{equation*}
\frac{\mathrm{d}b(\bar{w})}{\mathrm{d}\bar{w}}\Bigl|_{\,\bar{w}-\bar{w}_s\,=\,0.1}
  = -0.669,\ +0.316,\ +0.063
\qquad
\frac{\mathrm{d}b(\bar{w})}{\mathrm{d}\bar{w}}\Bigl|_{\,\bar{w}-\bar{w}_s\,=\,10}
  = -0.014,\ -0.001,\ +0.001
\end{equation*}}{formB_cmp_dbdw.png}

% ------------------------------------------------------------ step 3: a_bar'
\fpage{%
\begin{align*}
\bar{a}^{\prime}_{\text{A}} &= -\frac{\bigl(1-\bar{a}_{\mathrm{true}}\bigr)
   \ln\bigl(1-\bar{a}_{\mathrm{true}}\bigr)}{\bar{w}-\bar{w}_s}
&
\bar{a}^{\prime}_{\text{B}} &= \frac{\bar{a}_{\mathrm{true}}
   \bigl(1-\bar{a}_{\mathrm{true}}\bigr)}{\bar{w}-\bar{w}_s}
&
\bar{a}^{\prime}_{\text{C}} &= \frac{1-\bar{a}_{\mathrm{true}}}
   {1+\bar{w}-\bar{w}_s}
\end{align*}}{formB_cmp_aprime.png}

% -------------------------------------------------------- step 4: da_bar/db
\fpage{%
\begin{align*}
\frac{\partial\bar{a}_{\text{A}}}{\partial b}
  &= +\gp\bigl(1-\bar{a}_{\mathrm{true}}\bigr)
&
\frac{\partial\bar{a}_{\text{B}}}{\partial b}
  &= -\frac{\bar{a}_{\mathrm{true}}^{\,2}}{\bar{w}-\bar{w}_s}
&
\frac{\partial\bar{a}_{\text{C}}}{\partial b}
  &= -\frac{1}{1+\bar{w}-\bar{w}_s}
\end{align*}}{formB_cmp_dadb.png}

% ------------------------------------------------------------- step 5: J_b
\fpage{%
\begin{align*}
J_{b,\text{A}} &= \bigl(1-\bar{a}_{\mathrm{true}}\bigr)
   \Bigl[1+\ln\bigl(1-\bar{a}_{\mathrm{true}}\bigr)\Bigr]
&
J_{b,\text{B}} &= \frac{\bar{a}_{\mathrm{true}}^{\,2}
   \bigl(2\bar{a}_{\mathrm{true}}-1\bigr)}{\bigl(\bar{w}-\bar{w}_s\bigr)^{2}}
&
J_{b,\text{C}} &= \frac{1}{\gq^{2}}
\end{align*}

\begin{equation*}
J_b = 0:\quad
\text{A}:\ \bar{a}_{\mathrm{true}} = 1-e^{-1},\ \ \bar{w}-\bar{w}_s = 1.955
\qquad
\text{B}:\ \bar{a}_{\mathrm{true}} = \tfrac{1}{2},\ \ \bar{w}-\bar{w}_s = 1.128
\qquad
\text{C}:\ \text{none}
\end{equation*}}{formB_cmp_Jb.png}

% ---------------------------------------------- step 6: anchor vs b_eff
\fpage{%
\begin{align*}
\text{anchor}\quad
  b_0 &= \tfrac{9}{8}\ (\perp), \quad \tfrac{9}{16}\ (\parallel)
  &&\text{far-field reflection coefficient, one number, fixed} \\
\text{target}\quad
  b_{\mathrm{eff}}(\bar{w}) &= \text{page 1}
  &&\text{a curve: what the truth demands at each height}
\end{align*}

\begin{equation*}
\text{truth} = c_\perp\ \text{and}\ c_\parallel ;
\qquad
p = 1,\ \bar{w}_s = 1\ \text{pinned};
\qquad
\text{12 seeds}
\end{equation*}

\begin{center}
\begin{tabular}{llcccc}
\hline
truth & writing & $b_0$ & $b_{\mathrm{eff}}$ (hold) &
  $|b_0-b_{\mathrm{eff}}|$ & value of freeing it \\
\hline
$c_\perp$     & B & $1.1250$ & $1.1248$ & $0.0002$ & $+0.084$\ \ ($t=+1.63$) \\
$c_\parallel$ & C & $0.5625$ & $0.5571$ & $0.0054$ & $-0.016$\ \ ($t=-0.68$) \\
$c_\perp$     & C & $1.1250$ & $1.0587$ & $0.0663$ & $-1.610$\ \ ($t=-2.92$)\,$^\dagger$ \\
$c_\parallel$ & B & $0.5625$ & $0.3861$ & $0.1764$ & $-4.064$\ \ ($t=-6.48$)\,$^\dagger$ \\
\hline
\end{tabular}
\end{center}

\begin{equation*}
\text{value} := \bigl|\overline{e_a}\bigr|_{\text{free}}
              - \bigl|\overline{e_a}\bigr|_{\text{locked}}\ \ [\%],
\quad \text{final hold};
\qquad
\text{slope of the two large rows}: -24.3,\ -23.0
\end{equation*}

\emph{locked} pins the shape constant only (slot 5): its Kalman gain and
covariance are zero and it never leaves $b_0$. The gain state $\bar{a}$ is
estimated in both arms, so this is the marginal value of \emph{also} freeing
the shape constant. $^\dagger$ these two rows share a sign with the
$P(4,5)\!\to\!P_{44}$ inflation channel and await its control arm; the other
two do not.}{formB_anchor_vs_target.png}

% --------------------------------------------- step 7: gain error, perp
\fpage{%
\begin{equation*}
\text{truth} = c_\perp ,
\qquad
b_0 = \tfrac{9}{8}\ \text{for both writings},
\qquad
\sqrt{P[0]} = 0.0157\ (\text{B}),\ \ 0.0708\ (\text{C})
\end{equation*}

\begin{center}
\begin{tabular}{llccc}
\hline
writing & constant & descent peak \% & oscillation RMS \% & final hold \% \\
\hline
B & locked & $1.032 \pm 0.151$ & $0.341 \pm 0.075$ & $-0.482 \pm 0.368$ \\
B & free   & $1.102 \pm 0.219$ & $0.407 \pm 0.144$ & $-0.470 \pm 0.579$ \\
C & locked & $7.027 \pm 1.953$ & $2.762 \pm 1.205$ & $-5.579 \pm 2.328$ \\
C & free   & $6.605 \pm 2.652$ & $2.389 \pm 1.266$ & $-3.083 \pm 3.471$ \\
\hline
\end{tabular}
\end{center}

\begin{equation*}
\text{locked}\ \ \text{B}/\text{C} = 0.15,\ 0.12,\ 0.09
\qquad\Longrightarrow\qquad
\text{B ahead by}\ 6.8\text{--}11.6\times
\end{equation*}

The locked pair is the writing-versus-writing comparison with no estimator in
it. Page 1: $b_{\mathrm{eff},\text{B}}$ swings $0.0189$ here, $b_{\mathrm{eff},
\text{C}}$ swings $0.0699$.}{formB_page7_perp.png}

% --------------------------------------------- step 8: gain error, para
\fpage{%
\begin{equation*}
\text{truth} = c_\parallel\ \text{substituted into the wall-normal channel},
\qquad
b_0 = \tfrac{9}{16}\ \text{for both writings}
\end{equation*}

\begin{center}
\begin{tabular}{llccc}
\hline
writing & constant & descent peak \% & oscillation RMS \% & final hold \% \\
\hline
B & locked & $10.914 \pm 1.369$ & $4.572 \pm 1.177$ & $-9.238 \pm 1.789$ \\
B & free   & $\ \,9.728 \pm 2.624$ & $3.059 \pm 1.560$ & $-3.972 \pm 4.246$ \\
C & locked & $\ \,0.790 \pm 0.152$ & $0.375 \pm 0.138$ & $-0.506 \pm 0.201$ \\
C & free   & $\ \,0.783 \pm 0.190$ & $0.373 \pm 0.171$ & $-0.490 \pm 0.279$ \\
\hline
\end{tabular}
\end{center}

\begin{equation*}
\text{locked}\ \ \text{B}/\text{C} = 13.8,\ 12.2,\ 18.3
\qquad\Longrightarrow\qquad
\text{the order of page 7 is reversed}
\end{equation*}

$b_{\mathrm{eff}}$ swings $0.1769$ (B) against $0.0069$ (C) here, the mirror of
page 7. Form B pins a simple pole at $\bar{w}_s$; $c_\perp$ has exactly that
pole (residue $1$), $c_\parallel$ as implemented has none (finite $3.084$ at
contact). Two limits on reading this page: $c_\parallel$ is a substituted
truth, not the $x/y$ axes; and at \emph{equal} degrees of freedom --- Form B
with $b=\frac{9}{16}$, $p$ pinned and the origin free --- the shape residual is
$0.134\,\%$ against Form C's $0.404\,\%$, so the ordering above holds under the
rule that no constant may be fitted.}{formB_page8_para.png}

\end{document}
